\documentclass[10pt,a4paper]{amsart}
\usepackage{amsmath,amssymb,amsthm,graphicx}
% Retained minimal preamble for the Stacks Project chapter varieties.tex.
% Verbatim excerpt of the chapter preamble: Xy-pic setup (native preamble lines 14-18), the native theorem
% declarations with the mathematical macros (native preamble lines 154-217), and the project own no-op
% definitions of the reference/slogan/history display environments (native wrapper.tex lines 1-3) that the
% counted statement uses for its slogan. Copied mechanically; nothing added, removed or redefined.
% For commutative diagrams we use Xy-pic
\usepackage[all]{xy}

% We use 2cell for 2-commutative diagrams.
\xyoption{2cell}
% Theorem environments.
%
\theoremstyle{plain}
\newtheorem{theorem}[subsection]{Theorem}
\newtheorem{proposition}[subsection]{Proposition}
\newtheorem{lemma}[subsection]{Lemma}

\theoremstyle{definition}
\newtheorem{definition}[subsection]{Definition}
\newtheorem{example}[subsection]{Example}
\newtheorem{exercise}[subsection]{Exercise}
\newtheorem{situation}[subsection]{Situation}

\theoremstyle{remark}
\newtheorem{remark}[subsection]{Remark}
\newtheorem{remarks}[subsection]{Remarks}

\numberwithin{equation}{subsection}

% Macros
%
\def\lim{\mathop{\mathrm{lim}}\nolimits}
\def\colim{\mathop{\mathrm{colim}}\nolimits}
\def\Spec{\mathop{\mathrm{Spec}}}
\def\Hom{\mathop{\mathrm{Hom}}\nolimits}
\def\Ext{\mathop{\mathrm{Ext}}\nolimits}
\def\SheafHom{\mathop{\mathcal{H}\!\mathit{om}}\nolimits}
\def\SheafExt{\mathop{\mathcal{E}\!\mathit{xt}}\nolimits}
\def\Sch{\mathit{Sch}}
\def\Mor{\mathop{\mathrm{Mor}}\nolimits}
\def\Ob{\mathop{\mathrm{Ob}}\nolimits}
\def\Sh{\mathop{\mathit{Sh}}\nolimits}
\def\NL{\mathop{N\!L}\nolimits}
\def\CH{\mathop{\mathrm{CH}}\nolimits}
\def\proetale{{pro\text{-}\acute{e}tale}}
\def\etale{{\acute{e}tale}}
\def\QCoh{\mathit{QCoh}}
\def\Ker{\mathop{\mathrm{Ker}}}
\def\Im{\mathop{\mathrm{Im}}}
\def\Coker{\mathop{\mathrm{Coker}}}
\def\Coim{\mathop{\mathrm{Coim}}}

% Boxtimes
%
\DeclareMathSymbol{\boxtimes}{\mathbin}{AMSa}{"02}

%
% Macros for moduli stacks/spaces
%
\def\QCohstack{\mathcal{QC}\!\mathit{oh}}
\def\Cohstack{\mathcal{C}\!\mathit{oh}}
\def\Spacesstack{\mathcal{S}\!\mathit{paces}}
\def\Quotfunctor{\mathrm{Quot}}
\def\Hilbfunctor{\mathrm{Hilb}}
\def\Curvesstack{\mathcal{C}\!\mathit{urves}}
\def\Polarizedstack{\mathcal{P}\!\mathit{olarized}}
\def\Complexesstack{\mathcal{C}\!\mathit{omplexes}}
% \Pic is the operator that assigns to X its picard group, usage \Pic(X)
% \Picardstack_{X/B} denotes the Picard stack of X over B
% \Picardfunctor_{X/B} denotes the Picard functor of X over B
\def\Pic{\mathop{\mathrm{Pic}}\nolimits}
\def\Picardstack{\mathcal{P}\!\mathit{ic}}
\def\Picardfunctor{\mathrm{Pic}}
\def\Deformationcategory{\mathcal{D}\!\mathit{ef}}
\newenvironment{reference}{}{}
\newenvironment{slogan}{}{}
\newenvironment{history}{}{}
\usepackage{cleveref}
\crefname{originaldefinition}{original definition}{original definitions}
\Crefname{originaldefinition}{Original definition}{Original definitions}
\makeatletter
\newlabel{algebra-definition-geometrically-irreducible}{{037L}{source}{}{sourcepointer-1}{}}
\newlabel{algebra-definition-geometrically-irreducible@cref}{{[originaldefinition][0][]037L}{[1][0][]source}{}{}{}}
\makeatother
\crefname{originallemma}{original lemma}{original lemmas}
\Crefname{originallemma}{Original lemma}{Original lemmas}
\makeatletter
\newlabel{properties-lemma-characterize-irreducible}{{01OM}{source}{}{sourcepointer-2}{}}
\newlabel{properties-lemma-characterize-irreducible@cref}{{[originallemma][0][]01OM}{[1][0][]source}{}{}{}}
\makeatother
\title{Geometric irreducibility is Zariski local modulo connectedness: four equivalent formulations}
\author{The Stacks Project Authors}
\date{stacks-project:varieties.tex:local-snapshot-20261011; GNU FDL 1.2 or later}
\begin{document}
\maketitle
\noindent\textit{Editorial note.} All delivered mathematical text is the author's own, copied from the pinned Stacks Project chapter file varieties.tex; no statement or proof step was rewritten, reordered or removed.\par\medskip
\noindent\textit{Editorial note.} The statement carries the chapter's slogan environment, including the author's spelling of 'irreductibility'; it is reproduced verbatim and rendered by the no-op display environment of the retained preamble.\par\medskip
\noindent\textit{Editorial note.} The counted spans carry 2 navigation stubs and no bibliography citation.\par\medskip
\noindent\textit{Editorial note.} The preamble is the chapter's retained minimal preamble (sources/stacks-varieties/retained-preamble.tex lines 1-77), a verbatim excerpt of the native preamble.\par\medskip
\noindent\textit{Editorial note.} The label lemma-geometrically-irreducible-local is defined only in varieties.tex, so no cross-chapter restatement had to be weighed for this card.\par\medskip
\noindent\textit{Allowed standard prerequisites (editorial).}\begin{itemize}
\item Geometric irreducibility of a k-algebra. Algebra, Definition algebra-definition-geometrically-irreducible, invoked at proof lines 1413-1415 and 1429-1431. Supplies the affine translation used at the start of the proof.
\item Irreducibility of a scheme tested on affine opens. Properties, Lemma properties-lemma-characterize-irreducible, invoked at proof lines 1434-1436. Supplies that a scheme with irreducible affine opens is irreducible.
\item A cover by pairwise intersecting irreducible opens makes a space irreducible. Standard topology, used at proof lines 1439-1442. Supplies the last implication, applied after base change.
\end{itemize}\medskip
\ifdefined\hypertarget\hypertarget{sourcepointer-1}{}\fi
\noindent\textit{Source reference.} Original definition 037L: sources/stacks-chow/tags.txt, source lines 4204--4204. Pointer outside the counted span, invoked once for context; the load-bearing argument of the counted statement is the author's proof reproduced in full in the proof span.\par\medskip
\ifdefined\hypertarget\hypertarget{sourcepointer-2}{}\fi
\noindent\textit{Source reference.} Original lemma 01OM: sources/stacks-chow/tags.txt, source lines 2199--2199. Pointer outside the counted span, invoked once for context; the load-bearing argument of the counted statement is the author's proof reproduced in full in the proof span.\par\medskip
\begin{lemma}
\label{lemma-geometrically-irreducible-local}
\begin{slogan}
Geometric irreductibility is Zariski local modulo connectedness.
\end{slogan}
Let $k$ be a field. Let $X$ be a scheme over $k$.
The following are equivalent
\begin{enumerate}
\item $X$ is geometrically irreducible over $k$,
\item for every nonempty affine open $U$ the $k$-algebra $\mathcal{O}_X(U)$
is geometrically irreducible over $k$ (see
Algebra, Definition \ref{algebra-definition-geometrically-irreducible}),
\item $X$ is irreducible and there exists an affine open covering
$X = \bigcup U_i$ such that each $k$-algebra $\mathcal{O}_X(U_i)$ is
geometrically irreducible, and
\item there exists an open covering $X = \bigcup_{i \in I} X_i$
with $I \not = \emptyset$ such
that $X_i$ is geometrically irreducible for each $i$ and such that
$X_i \cap X_j \not = \emptyset$ for all $i, j \in I$.
\end{enumerate}
Moreover, if $X$ is geometrically irreducible so is every nonempty
open subscheme of $X$.
\end{lemma}

\begin{proof}
An affine scheme $\Spec(A)$ over $k$ is geometrically
irreducible if and only if $A$ is geometrically irreducible over $k$;
this is immediate from the definitions.
Recall that if a scheme is irreducible so is every nonempty
open subscheme of $X$, any two nonempty open subsets have
a nonempty intersection. Also, if every affine open is irreducible
then the scheme is irreducible, see Properties,
Lemma \ref{properties-lemma-characterize-irreducible}.
Hence the final statement of the lemma
is clear, as well as the implications (1) $\Rightarrow$ (2),
(2) $\Rightarrow$ (3), and (3) $\Rightarrow$ (4). If (4) holds,
then for any field extension $k'/k$ the scheme $X_{k'}$
has a covering by irreducible opens which pairwise intersect.
Hence $X_{k'}$ is irreducible. Hence (4) implies (1).
\end{proof}

\end{document}
